% Options for packages loaded elsewhere
\PassOptionsToPackage{unicode}{hyperref}
\PassOptionsToPackage{hyphens}{url}
\PassOptionsToPackage{dvipsnames,svgnames,x11names}{xcolor}
%
\documentclass[
  11pt,
]{article}
\usepackage{amsmath,amssymb}
\usepackage{iftex}
\ifPDFTeX
  \usepackage[T1]{fontenc}
  \usepackage[utf8]{inputenc}
  \usepackage{textcomp} % provide euro and other symbols
\else % if luatex or xetex
  \usepackage{unicode-math} % this also loads fontspec
  \defaultfontfeatures{Scale=MatchLowercase}
  \defaultfontfeatures[\rmfamily]{Ligatures=TeX,Scale=1}
\fi
\usepackage[]{charter}
\ifPDFTeX\else
  % xetex/luatex font selection
\fi
% Use upquote if available, for straight quotes in verbatim environments
\IfFileExists{upquote.sty}{\usepackage{upquote}}{}
\IfFileExists{microtype.sty}{% use microtype if available
  \usepackage[]{microtype}
  \UseMicrotypeSet[protrusion]{basicmath} % disable protrusion for tt fonts
}{}
\makeatletter
\@ifundefined{KOMAClassName}{% if non-KOMA class
  \IfFileExists{parskip.sty}{%
    \usepackage{parskip}
  }{% else
    \setlength{\parindent}{0pt}
    \setlength{\parskip}{6pt plus 2pt minus 1pt}}
}{% if KOMA class
  \KOMAoptions{parskip=half}}
\makeatother
\usepackage{xcolor}
\usepackage[margin=2.5cm]{geometry}
\usepackage{longtable,booktabs,array}
\usepackage{calc} % for calculating minipage widths
% Correct order of tables after \paragraph or \subparagraph
\usepackage{etoolbox}
\makeatletter
\patchcmd\longtable{\par}{\if@noskipsec\mbox{}\fi\par}{}{}
\makeatother
% Allow footnotes in longtable head/foot
\IfFileExists{footnotehyper.sty}{\usepackage{footnotehyper}}{\usepackage{footnote}}
\makesavenoteenv{longtable}
\setlength{\emergencystretch}{3em} % prevent overfull lines
\providecommand{\tightlist}{%
  \setlength{\itemsep}{0pt}\setlength{\parskip}{0pt}}
\setcounter{secnumdepth}{-\maxdimen} % remove section numbering
% ==========================================================================
%  IA na Sociedade — cabeçalho LaTeX (pdflatex)
%  Programa CiberExt 26-29 · FEELT38103
%  Adaptado de "AI in Society" (Universidade de Helsinque / Una Europa)
%  Licença: CC BY-NC-SA 4.0
% --------------------------------------------------------------------------
%  Espelha a identidade do ia-style.css: papel quente, tinta ardósia e a
%  faixa tricolor dos três eixos do curso (conceitos / aplicações / riscos).
% ==========================================================================

% ---------- Fontes ----------
% A serifada vem do build.sh, que passa  -V fontfamily=charter  ao pandoc.
% Aqui só acrescentamos a sem-serifa; ambas com fallback, porque nem toda
% instalação de LaTeX traz os mesmos pacotes.
\IfFileExists{helvet.sty}{\usepackage[scaled=0.90]{helvet}}{}
\usepackage{textcomp}
\IfFileExists{microtype.sty}{\usepackage{microtype}}{}

% ---------- Pacotes de layout ----------
\usepackage{xcolor}
\usepackage{tcolorbox}
\tcbuselibrary{skins,breakable}
\usepackage{titlesec}
\usepackage{fancyhdr}
\usepackage{enumitem}
\usepackage{amssymb}
\usepackage{booktabs}
\usepackage{array}
\usepackage{colortbl}
\usepackage{ragged2e}
\usepackage{fvextra}
\usepackage{newunicodechar}

% ---------- Caracteres Unicode que o pdflatex não conhece ----------
\newunicodechar{✓}{{\color{iaAplic}\checkmark}}
\newunicodechar{✅}{{\color{iaAplic}\checkmark}}
\newunicodechar{→}{\ensuremath{\rightarrow}}
\newunicodechar{←}{\ensuremath{\leftarrow}}
\newunicodechar{↔}{\ensuremath{\leftrightarrow}}
\newunicodechar{–}{--}
\newunicodechar{—}{---}
\newunicodechar{€}{\texteuro}
\newunicodechar{£}{\textsterling}
\newunicodechar{×}{\ensuremath{\times}}
\newunicodechar{≈}{\ensuremath{\approx}}
\newunicodechar{≥}{\ensuremath{\geq}}
\newunicodechar{≤}{\ensuremath{\leq}}
\newunicodechar{•}{\textbullet}
\newunicodechar{″}{''}
\newunicodechar{′}{'}

% ==========================================================================
%  PALETA — os mesmos valores do ia-style.css
% ==========================================================================
\definecolor{iaConceitos}{HTML}{24505C}   % petróleo — títulos, eixo 1
\definecolor{iaAplic}{HTML}{16746C}       % verdete  — subtítulos, eixo 2
\definecolor{iaRiscos}{HTML}{B45F14}      % açafrão  — destaques, eixo 3
\definecolor{iaAmeixa}{HTML}{77345F}      % links internos, termos técnicos
\definecolor{iaTinta}{HTML}{1F2933}       % corpo do texto
\definecolor{iaTintaFraca}{HTML}{5C6670}
\definecolor{iaRegra}{HTML}{E4DBCD}
\definecolor{iaBloco}{HTML}{F3EEE4}
\definecolor{iaNota}{HTML}{EAF1F0}
\definecolor{iaCaso}{HTML}{FCF3E6}
\definecolor{iaFilos}{HTML}{F7EFF4}

\color{iaTinta}

% ==========================================================================
%  ASSINATURA — faixa tricolor dos três eixos
% ==========================================================================
\newcommand{\faixaEixos}[1][\linewidth]{%
  \noindent
  \textcolor{iaConceitos}{\rule{0.3333#1}{4pt}}%
  \textcolor{iaAplic}{\rule{0.3333#1}{4pt}}%
  \textcolor{iaRiscos}{\rule{0.3333#1}{4pt}}%
  \par}

% ---------- Página de título ----------
\makeatletter
\newcommand{\subtitle}[1]{\gdef\@subtitle{#1}}
\providecommand{\@subtitle}{}
\def\@maketitle{%
  \newpage\null\vskip 1.5em%
  \faixaEixos
  \vskip 2.2em
  {\raggedright
    {\sffamily\bfseries\huge\color{iaConceitos}\@title\par}%
    \vskip 0.7em%
    {\sffamily\large\color{iaTintaFraca}\@subtitle\par}%
    \vskip 2em%
    {\color{iaRegra}\rule{\linewidth}{1pt}}%
    \vskip 1em%
    {\sffamily\small\color{iaTinta}\@author\par}%
    \vskip 0.5em%
    {\sffamily\footnotesize\color{iaTintaFraca}\@date\par}%
  }\par\vskip 2.5em}
\makeatother

% ==========================================================================
%  TÍTULOS DE SEÇÃO
% ==========================================================================
\titleformat{\section}[block]
  {\sffamily\Large\bfseries\color{iaConceitos}}
  {\thesection}{0.7em}{}
  [{\vspace{2pt}\color{iaRegra}\titlerule[1.2pt]}]
\titleformat{\subsection}
  {\sffamily\large\bfseries\color{iaAplic}}
  {\thesubsection}{0.6em}{}
\titleformat{\subsubsection}
  {\sffamily\normalsize\bfseries\color{iaRiscos}}
  {\thesubsubsection}{0.5em}{}
\titleformat{\paragraph}[runin]
  {\sffamily\small\bfseries\color{iaRiscos}}
  {}{0em}{}

\titlespacing*{\section}{0pt}{2.6ex plus 1ex minus .2ex}{1.3ex plus .2ex}
\titlespacing*{\subsection}{0pt}{2.2ex plus 1ex minus .2ex}{1ex plus .2ex}

% ==========================================================================
%  CABEÇALHO E RODAPÉ
% ==========================================================================
\setlength{\headheight}{24pt}
\renewcommand{\sectionmark}[1]{\markboth{#1}{}}
\pagestyle{fancy}
\fancyhf{}
\renewcommand{\headrulewidth}{0.6pt}
\renewcommand{\footrulewidth}{0pt}
\fancyhead[L]{\sffamily\footnotesize\color{iaAplic}Ética da IA}
\fancyhead[R]{\sffamily\footnotesize\color{iaAplic}\nouppercase{\leftmark}}
\fancyfoot[C]{\sffamily\footnotesize\color{iaConceitos}\thepage}
\renewcommand{\headrule}{%
  \hbox to\headwidth{\color{iaRegra}\leaders\hrule height \headrulewidth\hfill}}

% ==========================================================================
%  TEXTO CORRIDO
% ==========================================================================
\linespread{1.08}
\setlength{\parskip}{0.55em}
\setlength{\parindent}{0pt}
\emergencystretch=3em
\hbadness=10000

% ---------- Termos técnicos / código inline ----------
\let\oldtexttt\texttt
\renewcommand{\texttt}[1]{{\color{iaAmeixa}\small\ttfamily #1}}

% ---------- Blocos verbatim ----------
\fvset{breaklines=true,breakanywhere=true,fontsize=\small}

% ---------- Blocos de código (ambiente Shaded do pandoc) ----------
\renewenvironment{Shaded}
  {\begin{tcolorbox}[
      breakable, enhanced jigsaw, rounded corners,
      arc=3pt, boxrule=0pt, frame hidden,
      colback=iaBloco,
      borderline west={3pt}{0pt}{iaAmeixa},
      left=10pt, right=8pt, top=6pt, bottom=6pt,
      before skip=9pt, after skip=9pt]}
  {\end{tcolorbox}}

% ---------- Citações: caixa de caso ----------
\renewenvironment{quote}
  {\begin{tcolorbox}[
      breakable, enhanced jigsaw, rounded corners,
      arc=3pt, boxrule=0pt, frame hidden,
      colback=iaCaso,
      borderline west={3pt}{0pt}{iaRiscos},
      left=12pt, right=10pt, top=9pt, bottom=9pt,
      before skip=9pt, after skip=9pt]}
  {\end{tcolorbox}}

% ==========================================================================
%  BLOCOS DE APOIO
%  No markdown, escreva:   ::: nota   ...texto...   :::
%  O pandoc converte fenced_divs para os ambientes abaixo.
% ==========================================================================
\newtcolorbox{iabox}[3]{
  breakable, enhanced jigsaw, rounded corners,
  arc=3pt, boxrule=0pt, frame hidden,
  colback=#2,
  borderline west={3pt}{0pt}{#1},
  left=12pt, right=10pt, top=9pt, bottom=9pt,
  before skip=11pt, after skip=11pt,
  before upper={\setlength{\parskip}{0.5em}\setlength{\parindent}{0pt}},
  title={\sffamily\bfseries\footnotesize\textcolor{#1}{\MakeUppercase{#3}}},
  attach title to upper={\par\vspace{5pt}},
  fonttitle=\normalfont}

% Cada ambiente aceita um titulo opcional entre colchetes, vindo do
% atributo data-titulo do markdown:   \begin{filosofia}[Titulo] ... \end{filosofia}
% Os quatro primeiros espelham os icones do material original.
\newenvironment{filosofia}[1][Conceito filosofico]{\begin{iabox}{iaAmeixa}{iaFilos}{#1}}{\end{iabox}}
\newenvironment{tecnica}[1][Nota tecnica]{\begin{iabox}{iaConceitos}{iaNota}{#1}}{\end{iabox}}
\newenvironment{contexto}[1][Contexto]{\begin{iabox}{iaAplic}{iaBloco}{#1}}{\end{iabox}}
\newenvironment{definicao}[1][Definicao]{\begin{iabox}{iaRiscos}{iaCaso}{#1}}{\end{iabox}}

\newenvironment{nota}[1][Nota]{\begin{iabox}{iaConceitos}{iaNota}{#1}}{\end{iabox}}
\newenvironment{caso}[1][Caso]{\begin{iabox}{iaRiscos}{iaCaso}{#1}}{\end{iabox}}
\newenvironment{reflexao}[1][Para pensar]{\begin{iabox}{iaAplic}{iaBloco}{#1}}{\end{iabox}}
\newenvironment{licenca}
  {\par\vspace{2em}\faixaEixos\par\vspace{0.8em}
   \begingroup\sffamily\scriptsize\color{iaTintaFraca}}
  {\endgroup\par}

% ==========================================================================
%  LISTAS
% ==========================================================================
\setlist[itemize]{leftmargin=1.4em, itemsep=2pt, topsep=4pt, parsep=0pt}
\setlist[enumerate]{leftmargin=1.7em, itemsep=2pt, topsep=4pt, parsep=0pt}
\renewcommand{\labelitemi}{\textcolor{iaRiscos}{\textbullet}}
\renewcommand{\labelitemii}{\textcolor{iaAplic}{--}}

% ==========================================================================
%  TABELAS
% ==========================================================================
\renewcommand{\arraystretch}{1.3}
\arrayrulecolor{iaRegra}

% ==========================================================================
%  RÓTULOS EM PORTUGUÊS
%  Definidos aqui (e não via -V na linha de comando) porque acentos passados
%  como argumento dependem do locale do terminal e podem se corromper.
% ==========================================================================
\renewcommand{\contentsname}{Sumário}
\renewcommand{\refname}{Referências}
\renewcommand{\figurename}{Figura}
\renewcommand{\tablename}{Tabela}
\ifLuaTeX
  \usepackage{selnolig}  % disable illegal ligatures
\fi
\IfFileExists{bookmark.sty}{\usepackage{bookmark}}{\usepackage{hyperref}}
\IfFileExists{xurl.sty}{\usepackage{xurl}}{} % add URL line breaks if available
\urlstyle{same}
\hypersetup{
  pdftitle={Ética da Inteligência Artificial},
  pdfauthor={Programa CiberExt 26-29 · FEELT38103 · uso do docente},
  colorlinks=true,
  linkcolor={iaAmeixa},
  filecolor={Maroon},
  citecolor={Blue},
  urlcolor={iaRiscos},
  pdfcreator={LaTeX via pandoc}}

\title{Ética da Inteligência Artificial}
\usepackage{etoolbox}
\makeatletter
\providecommand{\subtitle}[1]{% add subtitle to \maketitle
  \apptocmd{\@title}{\par {\large #1 \par}}{}{}
}
\makeatother
\subtitle{Capítulo 1 --- Gabarito e critérios}
\author{Programa CiberExt 26-29 · FEELT38103 · uso do docente}
\date{2026}

\begin{document}
\maketitle

{
\hypersetup{linkcolor=iaConceitos}
\setcounter{tocdepth}{2}
\tableofcontents
}
\hypertarget{capuxedtulo-1-gabarito-e-crituxe9rios}{%
\section{Capítulo 1 --- Gabarito e
critérios}\label{capuxedtulo-1-gabarito-e-crituxe9rios}}

\begin{nota}[Material de uso do docente]

Não publicar junto com o curso. Existe no repositório porque faz parte
da entrega acadêmica do projeto de extensão.

\end{nota}

\hypertarget{exercuxedcios-1a-e-1b-valores-e-normas}{%
\subsection{Exercícios 1a e 1b --- Valores e
normas}\label{exercuxedcios-1a-e-1b-valores-e-normas}}

\textbf{Não há resposta certa.} São instrumentos de autoconhecimento,
adaptados da escala de valores de Schwartz (2001). Todos os
participantes recebem o ponto por responderem.

\textbf{O que observar na discussão em fórum:} o interesse pedagógico
está na diferença entre as duas listas. Valores que a pessoa marca como
pessoais no 1a com frequência não são os que ela prioriza no 1b, quando
age como formuladora de política pública. Vale pedir que os estudantes
comparem as próprias respostas --- essa dissonância é o ponto do
exercício e prepara o terreno para a guilhotina de Hume.

\begin{center}\rule{0.5\linewidth}{0.5pt}\end{center}

\hypertarget{exercuxedcio-2-associauxe7uxe3o-valores-normas}{%
\subsection{Exercício 2 --- Associação valores ×
normas}\label{exercuxedcio-2-associauxe7uxe3o-valores-normas}}

\begin{longtable}[]{@{}
  >{\raggedright\arraybackslash}p{(\columnwidth - 2\tabcolsep) * \real{0.5000}}
  >{\raggedright\arraybackslash}p{(\columnwidth - 2\tabcolsep) * \real{0.5000}}@{}}
\toprule\noalign{}
\begin{minipage}[b]{\linewidth}\raggedright
Norma
\end{minipage} & \begin{minipage}[b]{\linewidth}\raggedright
Valor
\end{minipage} \\
\midrule\noalign{}
\endhead
\bottomrule\noalign{}
\endlastfoot
1. Proteções adicionais sobre dados de treinamento & Privacidade \\
2. Promover justiça e eliminar discriminação & Igualdade e não
discriminação \\
3. Libertar as pessoas do trabalho rotineiro & Autorrealização e
florescimento humano \\
4. Não manipular cidadãos em eleições & Democracia \\
5. Reduzir a pegada de carbono da IA & Proteção da natureza \\
6. Não prejudicar a liberdade de estabelecer padrões próprios &
Autonomia \\
7. Robustez técnica e proteção contra uso malicioso & Segurança \\
\end{longtable}

\textbf{Justificativa geral:} cada norma é a formulação prescritiva
(``deve-se\ldots{}'') de um valor subjacente. O exercício treina
exatamente a distinção valor/norma da seção III do capítulo.

\textbf{Confusões esperadas:} os itens 1 e 7 se confundem (privacidade ×
segurança) --- a diferença é que o item 1 trata de controle sobre os
próprios dados, e o 7, de resistência a ataques. Os itens 3 e 6 também
se confundem (autorrealização × autonomia) --- o 3 é sobre poder
desenvolver-se, o 6, sobre não ter o próprio juízo substituído.

\begin{center}\rule{0.5\linewidth}{0.5pt}\end{center}

\hypertarget{exercuxedcio-3-tuitando-sobre-questuxf5es-uxe9ticas}{%
\subsection{Exercício 3 --- Tuitando sobre questões
éticas}\label{exercuxedcio-3-tuitando-sobre-questuxf5es-uxe9ticas}}

\textbf{Sem resposta única.} Avaliado por revisão por pares.

Uma resposta adequada deve:

\begin{itemize}
\tightlist
\item
  tomar uma posição explícita sobre o caso, e não apenas descrevê-lo;
\item
  reconhecer que o comportamento do modelo reflete padrões estatísticos
  dos dados de treinamento, não uma ``opinião'' da máquina;
\item
  evitar as duas saídas fáceis: ``é só um algoritmo, não tem culpa'' e
  ``a tecnologia é racista/sexista por natureza''.
\end{itemize}

\textbf{Erros comuns:} confundir a explicação técnica (por que o modelo
faz isso) com a justificação ética (se isso é aceitável) --- é a
guilhotina de Hume em ação; tratar o caso como problema puramente
técnico, solucionável com mais dados.

\begin{center}\rule{0.5\linewidth}{0.5pt}\end{center}

\hypertarget{exercuxedcio-3c-ensaio-150-a-300-palavras}{%
\subsection{Exercício 3c --- Ensaio (150 a 300
palavras)}\label{exercuxedcio-3c-ensaio-150-a-300-palavras}}

Uma resposta adequada deve identificar ao menos três dos seguintes
elementos:

\begin{itemize}
\tightlist
\item
  \textbf{Valores em jogo:} igualdade, dignidade, representação justa.
\item
  \textbf{Normas:} a IA não deve reforçar estereótipos; sistemas devem
  ser testados quanto a viés antes da implantação.
\item
  \textbf{Princípios:} equidade e não discriminação (cap. 6),
  transparência (cap. 4), responsabilização (cap. 3).
\item
  \textbf{A distinção fato/valor:} o fato de os dados serem enviesados
  não determina, sozinho, que seja errado --- é preciso um argumento
  normativo adicional.
\end{itemize}

\begin{longtable}[]{@{}
  >{\raggedright\arraybackslash}p{(\columnwidth - 6\tabcolsep) * \real{0.2500}}
  >{\raggedright\arraybackslash}p{(\columnwidth - 6\tabcolsep) * \real{0.2500}}
  >{\raggedright\arraybackslash}p{(\columnwidth - 6\tabcolsep) * \real{0.2500}}
  >{\raggedright\arraybackslash}p{(\columnwidth - 6\tabcolsep) * \real{0.2500}}@{}}
\toprule\noalign{}
\begin{minipage}[b]{\linewidth}\raggedright
Critério
\end{minipage} & \begin{minipage}[b]{\linewidth}\raggedright
Insuficiente
\end{minipage} & \begin{minipage}[b]{\linewidth}\raggedright
Adequado
\end{minipage} & \begin{minipage}[b]{\linewidth}\raggedright
Excelente
\end{minipage} \\
\midrule\noalign{}
\endhead
\bottomrule\noalign{}
\endlastfoot
Identifica valores e normas & não distingue os dois & distingue e nomeia
& articula a relação entre eles \\
Julga a importância & apenas descreve & julga com justificativa &
pondera argumentos contrários \\
Usa conceitos do capítulo & ausentes & presentes & aplicados com
precisão \\
Extensão e clareza & fora do limite & dentro do limite & bem
estruturado \\
\end{longtable}

\begin{center}\rule{0.5\linewidth}{0.5pt}\end{center}

\begin{licenca}

Gabarito elaborado a partir de \emph{Ethics of AI}, Universidade de
Helsinque. Tradução e adaptação: José Lucas Lira Bizil, Fernando Mazzeto
Lisboa Lima e Matheus da Silva Fernandes, Programa CiberExt 26-29
(FEELT38103), UFU, 2026. Licença CC BY-NC-SA 4.0.

\end{licenca}

\end{document}
